\chapter{CONSIDERAÇÕES FINAIS}
\label{cap:conclusao}
No que tange ao sistema de controle e estabilização, os ensaios práticos evidenciaram que a abordagem LQR/LQG, em conjunto com o Filtro de Kalman linear, manteve a estabilidade de rolamento e arfagem do veículo. Em ensaios com o deslocamento horizontal restringido por um apoio físico, a altura de referência foi alcançada; em voo livre, porém, o veículo não manteve a posição nos eixos $X$ e $Y$ nem a guinada durante a subida. Essa deriva o expôs a forças externas do ambiente de ensaio, como cabos, paredes e correntes de ar, que provocaram a inclinação excessiva e o acionamento do sistema de segurança (\textit{failsafe}). A manutenção da posição horizontal mostrou-se, assim, o principal ponto a aperfeiçoar, pois expôs o impacto de perturbações não modeladas e de torques constantes gerados por assimetrias físicas.

Diante desses resultados, como trabalhos futuros na área de controle, sugere-se o aprimoramento das técnicas de modelagem matemática da planta, buscando uma representação que englobe efeitos aerodinâmicos e aproxime o modelo no espaço de estados do comportamento físico real. Recomenda-se também um rigor maior no projeto mecânico para garantir o balanceamento adequado da carga estrutural no drone, alinhando o centro de gravidade de modo a mitigar os torques parasitas. Por fim, propõe-se a adoção de técnicas de estimação de estados mais sofisticadas, substituindo o filtro linear pelo Filtro de Kalman Estendido (EKF -- \textit{Extended Kalman Filter}) baseado na cinemática de quatérnios. Em contraste com os ângulos de Euler tradicionais, o uso de quatérnios eliminaria as singularidades matemáticas (\textit{gimbal lock}) e reduziria substancialmente a carga de cálculos trigonométricos no microcontrolador, proporcionando uma fusão sensorial mais robusta, computacionalmente eficiente e adequada à dinâmica inerentemente não linear do VANT.
% Comunicação Raspberry Pi <-> STM32
% ============================================================================
% CONSIDERAÇÕES FINAIS — comunicação entre o Raspberry Pi e o STM32
% Chamado em partes/textuais/conclusao.tex.
% ============================================================================


Na comunicação entre o computador de missão e o controlador de voo, foi desenvolvido o DroneLink, protocolo de aplicação sobre o enquadramento MAVLink 2, definido em um único arquivo XML do qual são gerados os códigos do STM32 e do Raspberry Pi. O protocolo divide a missão em trechos executados um de cada vez: cada trecho é validado pelo STM32 antes de ser aceito, só é dado como concluído quando o VANT está estável no alvo e pode ser retransmitido com segurança, pois pedidos repetidos não são executados duas vezes. A identidade de \textit{boot} persistente impede a execução de comandos anteriores a uma reinicialização, e a supervisão do \textit{HEARTBEAT} e da autorização do operador leva o VANT a manter a posição e, em seguida, a pousar se o Raspberry Pi deixar de responder.

Sobre esse protocolo, o Raspberry Pi oferece ao operador comandos com a nomenclatura aeronáutica usual, entre eles o \texttt{route}, que ordena os pontos pelo algoritmo genético e executa a rota de forma automática. No ensaio integrado pela UART, com o VANT simulado no STM32, não foram registradas perdas nem erros de \textit{checksum} durante os 60~s de observação do enlace, com latência de ida e volta em torno de 10~ms; todas as recusas ocorreram com o motivo previsto; um pedido repetido foi executado uma única vez; a missão completa foi executada com erro máximo de 1,4~cm por ponto; e o STM32 manteve a posição em cerca de 1~s e iniciou o pouso entre 5,29 e 5,98~s diante da perda da autorização ou do encerramento do programa do Raspberry Pi. Os ensaios também permitiram identificar e corrigir seis falhas de integração, três delas visíveis somente com os dois lados conectados.

Como trabalhos futuros, recomenda-se realizar ensaios integrados com os motores em funcionamento, avaliando a comunicação sob vibração e interferência elétrica; reduzir o tempo de cálculo do algoritmo genético no Raspberry Pi Zero W; e avaliar a assinatura de mensagens do MAVLink 2 para autenticação da origem, especialmente caso a comunicação seja disponibilizada por rede sem fio.


% Planejamento de rotas
% ============================================================================
% CONSIDERAÇÕES FINAIS — planejamento de rotas
% Chamado em partes/textuais/conclusao.tex.
% ============================================================================

No planejamento de rotas, a ordem de visita dos pontos foi modelada como o Problema do Caixeiro Viajante com custo igual à distância tridimensional do circuito, que, em um VANT, relaciona-se diretamente com o tempo de voo e o consumo de energia. A comparação entre a ordem aleatória, o algoritmo guloso, o algoritmo genético (AG) e dois métodos exatos (a força bruta e o Held-Karp) mostrou que a obtenção do ótimo global deixa de ser viável com poucas dezenas de pontos. A força bruta foi utilizada somente até 10 pontos, pois com 11 a instância já levou 18~s e, com 20, levaria cerca de 32 mil anos. O Held-Karp, que obtém o mesmo ótimo por programação dinâmica, foi utilizado até 22 pontos, com cerca de 1~min por instância; com 23, ultrapassou 2~min e exigiu 1,7~GB, sendo retirado do comparativo, e a partir de 25 a memória necessária (7,5~GB) excedeu a do computador utilizado. Assim, a partir de 23 pontos, a comparação restringiu-se ao AG, ao guloso e à ordem aleatória.

O algoritmo genético mostrou-se viável em toda a faixa avaliada, de 3 a 30 pontos, com tempo de cálculo entre 0,13 e 0,45~s no computador pessoal, alocação de memória inferior a 0,1~MB e rotas com qualidade média entre 96,5\% e 100\% da referência, tendo obtido o ótimo global em 150 das 200 instâncias em que ele era conhecido. O algoritmo guloso, embora mais rápido, produziu rotas de 6\% a 18\% mais longas em média. Por ter sido validado nas faixas em que o ótimo é conhecido, o AG foi adotado também nas instâncias maiores, nas quais o ótimo não pode ser calculado e suas rotas se mantiveram próximas da melhor rota conhecida. A execução de rotas tridimensionais pelo controlador de voo, com o VANT simulado no STM32, confirmou que a ordem calculada é executada conforme planejada, com erro máximo de 3,1~cm em cada ponto e redução de 21,3\% na distância total da rota do AG em relação à do guloso.

O mesmo AG, executado no Raspberry Pi Zero W, produziu rotas idênticas às do computador pessoal, verificadas por assinatura SHA-256, em um intervalo de 10 a 24~s, tempo considerado aceitável para um cálculo pré-voo. O algoritmo genético constituiu, assim, uma solução viável, adotada e validada no protótipo, abrangendo sua execução no microprocessador embarcado e a integração com o controlador de voo. Essa validação corrobora sua viabilidade operacional.

Como trabalhos futuros para esta frente, recomenda-se reduzir o tempo de cálculo no Raspberry Pi, seja pela diminuição do número de gerações, seja pela adoção de uma implementação de código compilado; aprofundar o estudo da busca local \textit{2-opt} com múltiplas partidas, que, na investigação exploratória, obteve rotas superiores às do AG em cerca de 1/30 do tempo, seja como alternativa exclusiva, seja combinada ao AG por meio de um algoritmo memético; e incorporar à função de custo variáveis e restrições dinâmicas de tempo, energia e obstáculos, características inerentes a operações logísticas e missões reais.

